\documentclass[10pt,a4paper]{amsart}
\usepackage[margin=19mm]{geometry}
\usepackage{amsmath,amssymb,mathtools}
\usepackage[scr=rsfs,cal=euler]{mathalfa}
\usepackage{xfrac}
\usepackage[unicode,hidelinks,bookmarks=false]{hyperref}
\newcommand{\ie}{i.e.\ }
\newcommand{\cf}{cf.\ }
\newcommand{\resp}{respectively\ }
\newcommand{\Subsection}{\subsection}
\providecommand{\nobreakdash}{\nobreak\hbox{-}}
\setlength{\emergencystretch}{2em}
\multlinegap0pt
\usepackage{siunitx}
\usepackage{siunitx}
\usepackage{siunitx}
\usepackage{siunitx}
\usepackage[shortlabels]{enumitem}
\usepackage{mathtools}
\usepackage{siunitx}
\usepackage{algorithm}\usepackage{algpseudocode}
\newtheorem{theorem}{Theorem}
\title{Expand More, Shrink Less: Shaping Effective-Rank Dynamics for Dense Scaling in Recommendation}
\author{Guoming Li, Shangyu Zhang, Junwei Pan, Wentao Ning, Jin Chen, Gengsheng Xue, Chao Zhou, Shudong Huang, Haijie Gu, Menglin Yang}
\date{Source: arXiv:2605.23191v1 (CC BY 4.0)}
\begin{document}
\maketitle

\noindent\emph{Source and licence.} Expand More, Shrink Less: Shaping Effective-Rank Dynamics for Dense Scaling in Recommendation. Version arXiv:2605.23191v1 (CC BY 4.0), distributed by arXiv under the Creative Commons Attribution 4.0 International licence, http://creativecommons.org/licenses/by/4.0/, as recorded in the arXiv licence metadata for this exact version and shown on the abstract page https://arxiv.org/abs/2605.23191v1. Full licence notice: \texttt{licenses/arxiv-2605.23191-CC-BY-4.0.txt}.\par\smallskip

\noindent\textit{Editorial scope.} Counted unit: the article's theorem thm:expressivity (source0.tex lines 1544--1558). The article's complete original proof of it is delivered with it (source0.tex lines 1560--1571, one \texttt{proof} environment of 12 lines). No mathematics is written, repaired or re-derived: the statement and the proof are the article's own lines, in the article's own notation, and no retained line is substituted or re-flowed.  No further source-local context is needed: every cross-reference of every retained line resolves inside the two delivered spans.  Nothing was omitted from any retained span, so the package carries no omission record.  No retained line contains a citation, so the wrapper carries no bibliography and no bibliography entry.  No retained line contains a figure environment, an image-inclusion command or drawing code, so no image is delivered and the package declares no asset dependency.  Editorial formatting only, in addition: an AMS \texttt{amsart} preamble that restates the article's own declarations and macros used by the retained lines, namely the environments \texttt{theorem} on separate counters, together with the standard packages those lines need.  The retained environments print in this document as Theorem 1 in order of appearance, whereas the article prints the same items with the numbering of its own full document.  The complete native source of the article as distributed in the arXiv e-print for this exact version is bundled unchanged as \texttt{sources/201192/source0.tex}.
\begin{theorem}[Fine-grained expressivity and subspace constraint]
\label{thm:expressivity}
Define the reachable set
\[
\mathcal R_{d^{\ast}}(X)
=\bigl\{\,C\in\mathbb R^{T\times D}\mid \exists\,\mathbf W,\ 
C=\Phi_{d^{\ast}}(X;\mathbf W)\,\bigr\}.
\]
Assume \(X\neq 0\). Then:
\begin{itemize}[leftmargin=15pt,parsep=2pt,itemsep=2pt,topsep=2pt]
  \item[(i)] (\textbf{Universal reachability at \(d^{\ast}=1\)}) \(\mathcal R_1(X)=\mathbb R^{T\times D}\).
  \item[(ii)] (\textbf{Strict deficiency at \(d^{\ast}>1\)}) For every \(d^{\ast}>1\), \(\mathcal R_{d^{\ast}}(X)\subsetneq\mathcal R_1(X)\).
  \item[(iii)] (\textbf{Dependence on effective rank}) Let \(V_{\mathrm{blocks}}=\operatorname{span}\{\mathbf b_1,\dots,\mathbf b_K\}\subset\mathbb R^{d^{\ast}}\). All perturbations produced by \(\Phi_{d^{\ast}}\) have each block lying in \(V_{\mathrm{blocks}}\). If \(\operatorname{erank}(X)\ll\min\{T,D\}\), then the \(\mathbf b_k\) are approximately confined to a low-dimensional subspace of \(\mathbb R^{d^{\ast}}\), so there exist fine-grained variations (directions orthogonal to \(V_{\mathrm{blocks}}\)) that \(\Phi_{d^{\ast}}\) cannot express.
\end{itemize}
\end{theorem}

\begin{proof}
We structure the proof into two parts: (i) the case $d^{\ast}=1$, and (ii)–(iii) the case $d^{\ast}>1$. 
\subsubsection*{\bf Part (i).} For $d^{\ast}=1$, when \(d^{\ast}=1\) each block is a scalar and \(K=N\). The update term becomes \(\mathbf W\mathbf x\) and \(\operatorname{vec}(C)=\mathbf x+\mathbf W\mathbf x=(\mathbf I+\mathbf W)\mathbf x\). Since \(\mathbf x\neq 0\), for any target vector \(\mathbf z\in\mathbb R^N\) we can choose \(\mathbf W\) (e.g. via the outer-product construction \(\mathbf W=(\mathbf z-\mathbf x)\mathbf x^\top/(\mathbf x^\top\mathbf x)\)) so that \((\mathbf I+\mathbf W)\mathbf x=\mathbf z\). Hence \(\mathcal R_1(X)=\mathbb R^{T\times D}\).

\subsubsection*{\bf Part (ii) and (iii).} For general \(d^{\ast}>1\), write the perturbation \(\mathbf p=\operatorname{vec}(C-X)=(\mathbf W\otimes\mathbf I_{d^{\ast}})\mathbf x\) and partition \(\mathbf p\) into blocks \(\mathbf p_1,\dots,\mathbf p_K\). By block-Kronecker structure the \(k\)-th block satisfies
\[
\mathbf p_k=\sum_{j=1}^K W_{kj}\,\mathbf b_j,
\]
hence \(\mathbf p_k\in V_{\mathrm{blocks}}\) for every \(k\). Thus every reachable perturbation is blockwise constrained to the subspace \(V_{\mathrm{blocks}}\subseteq\mathbb R^{d^{\ast}}\); equivalently, \(\mathcal R_{d^{\ast}}(X)\) lies inside a linear submanifold of \(\mathbb R^{T\times D}\) of strictly smaller dimension whenever \(\dim(V_{\mathrm{blocks}})<d^{\ast}\). Therefore, whenever \(V_{\mathrm{blocks}}\neq\mathbb R^{d^{\ast}}\) there exist target matrices whose required block perturbation contains a component orthogonal to \(V_{\mathrm{blocks}}\); such targets are not in \(\mathcal R_{d^{\ast}}(X)\), so \(\mathcal R_{d^{\ast}}(X)\subsetneq\mathcal R_1(X)\).

On the root of low effective rank makes the deficiency generic, note that \(\operatorname{erank}(X)=\|X\|_F^2/\|X\|_2^2\ll\min\{T,D\}\) implies that most energy of \(X\) is concentrated in a few global singular directions. Under standard vectorization (row- or column-major) contiguous local blocks \(\mathbf b_k\) therefore inherit strong correlations and typically span a low-dimensional subspace of \(\mathbb R^{d^{\ast}}\). Concretely, \(\dim(V_{\mathrm{blocks}})\ll d^{\ast}\) in this regime, so many fine-grained directions orthogonal to \(V_{\mathrm{blocks}}\) exist and cannot be generated by the Kronecker-constrained operator. This establishes the dependence on effective rank and completes the proof.
\end{proof}

\end{document}
